\documentclass[book.tex]{subfiles}
\begin{document}

\chapter{OpenKIM and ColabFit}
\label{chap:hamiltonian_nn}
\noindent\rule{11cm}{0.4pt} \\
\textbf{Prerequisites:} \autoref{chap:hamiltonians_mechanics}, \autoref{chap:neural_ode} \\
\textbf{Difficulty Level:} ** \\
\noindent\rule{11cm}{0.4pt}
\vspace{5mm}



\section{Introduction}

New computers have enabled more accurate predictions with more rigorous models for molecular simulations. Post HF to DFT to Multi-body potentials to Pair potentials provide a ladder of approximations.

There are many possibilities. For example, over 60 interatomic potentials have been designed for silicon alone. For example, the Stillinger-Weber silicon approximation predicts silicon is ductile while it is actually brittle. Choice of interatomic potential is important.

The major challenge for ML potential development is data. OpenKim is online infrastructure for testing and deploying interatomic potentials. Colabfit is a database of data for interatomic potential fitting. Many databases only store ground state energies, while we need access to non-equilibrium energies to train non-equilibrium energies. See at \url{colabfit.org}. The database includes 50 datasets. Builds on the OpenCatalyst database. All of this is stored in MongoDB.

KLIFF (Kim-based Learning Integrated Fitting Framework) is a new framework that finds a niche between DeepMDKit and Torch/Jax MD. It supports deployment to LAMMPS and multi-GPU training. It also support uncertainty quantification techniques. Includes LibDescriptor, a differentiable descriptor library written in C++. OpenKim \url{openkim.org} has been running since 2009. Has a chat-system ChatKIM to help the user select an appropriate recommender system.


\section{Couple Electrochemo-mechanical Batter Models}






\end{document}